% RamanujanDreams-CUDA: High-Performance CUDA Implementation of
% Conservative Matrix Field Walks for pFq Hypergeometric Functions
%
% Benchmark paper / preprint for EuroHPC project ehpc916

\documentclass[11pt,a4paper]{article}
\usepackage[utf8]{inputenc}
\usepackage{amsmath,amssymb,amsthm}
\usepackage{booktabs}
\usepackage{graphicx}
\usepackage{hyperref}
\usepackage{listings}
\usepackage{xcolor}
\usepackage{geometry}
\usepackage{siunitx}
\usepackage{multirow}

\geometry{margin=2.5cm}

\title{RamanujanDreams-CUDA: GPU-Accelerated Conservative Matrix Field Walks \\
       for pFq Hypergeometric Functions on NVIDIA H100}
\author{David Vesterlund \\
        \small EuroHPC project ehpc916, MareNostrum5 (BSC)}
\date{\today}

\begin{document}
\maketitle

\begin{abstract}
We present three CUDA implementations of the Conservative Matrix Field (CMF) walk
algorithm for generalized hypergeometric functions ${}_pF_q$, targeting NVIDIA
H100 GPUs on the MareNostrum5 supercomputer. The implementations use three
distinct arithmetic strategies: native double precision (f64), Residue Number
System (RNS) exact modular arithmetic, and double-single (df64) emulated
extended precision. We benchmark all three variants across five CMF types
(${}_3F_2$, ${}_4F_3$, ${}_5F_4$, ${}_6F_5$, ${}_8F_7$) with 10{,}000
trajectories each at walk depth $N=1000$, using both random initial points
(limit search) and the classical zeta initial points
$(1,\ldots,1;\,2,\ldots,2)$ with $z=1$ (delta search).

The implementations are validated against the reference \texttt{ramanujantools}
Python library: matrix construction, walk ordering, $p/q$ extraction, and the
delta (irrationality measure) formula are verified to match. A key technical
contribution is log-space scale tracking through projective normalization,
enabling correct delta computation
$\delta = -\bigl(1 + \log|L - p/q| / \log q\bigr)$
despite per-step rescaling of the walk matrix.

Results show that df64 achieves the highest throughput for limit search
(up to \SI{789000}{traj/s} for ${}_3F_2$), while f64 is within 5--20\%.
RNS provides exact arithmetic at $\sim$100$\times$ lower throughput but
with guaranteed correctness, making it suitable for verification and deep
walks. For the standard Ramanujan Machine workflow of GPU-based limit
sweeping followed by CPU-based delta refinement, the df64 and f64 variants
provide a \textbf{100--1000$\times$ speedup} over single-threaded Python,
enabling real-time exploration of initial point spaces.
\end{abstract}

\section{Introduction}

The Ramanujan Machine project \cite{ramanujan-machine} discovers mathematical
constants through algorithmic exploration of Conservative Matrix Fields (CMFs)
and Polynomial Continued Fractions (PCFs). A CMF defines a family of matrices
satisfying a conservation property, enabling the computation of limits and
irrationality measures (deltas) for generalized hypergeometric functions
${}_pF_q$.

The standard workflow has two phases:
\begin{enumerate}
    \item \textbf{Limit search}: Sweep many initial points to discover
    promising limit values. This is embarrassingly parallel and well-suited
    for GPU acceleration.
    \item \textbf{Delta refinement}: For a known good initial point, search
    over different trajectories to find the best irrationality measure.
    This is traditionally done on CPU due to the need for exact arithmetic.
\end{enumerate}

This paper presents GPU-accelerated implementations of the CMF walk using
three arithmetic strategies, benchmarked on NVIDIA H100 GPUs, and investigates
whether GPU-based delta search is viable.

\section{Background}

\subsection{Conservative Matrix Fields}

A CMF is defined by a set of axis matrices $M_{x_i}$ and $M_{y_j}$ satisfying
the conservation property:
\[
M_x(x, y) \cdot M_y(x+1, y) = M_y(x, y) \cdot M_x(x, y+1)
\]
For the ${}_pF_q$ CMF, the axis matrix is:
\[
M_{\text{axis}} = I + \frac{C}{a}
\]
where $C$ is the companion matrix of the hypergeometric differential equation:
\[
D(\theta) = \theta \prod_{i=0}^{q-1}(\theta + y_i - 1) - z \prod_{j=0}^{p-1}(\theta + x_j)
\]
The rank is $\max(p, q+1)$, reduced by one when $z=1$ and $p = q+1$.

\subsection{The Matrix Walk}

Given a trajectory direction $\mathbf{d}$ and starting position $\mathbf{s}$,
the walk matrix at depth $N$ is:
\[
W_N = T(0) \cdot T(1) \cdots T(N-1)
\]
where $T(n)$ is the trajectory matrix at step $n$. The limit is extracted as
the ratio $p/q$ from the walk matrix using the default initial and final
projection vectors:
\[
p = (e_0 \cdot W_N \cdot e_{-1}), \qquad q = (e_1 \cdot W_N \cdot e_{-1})
\]
matching the \texttt{ramanujantools} convention.

\subsection{Irrationality Measure (Delta)}

The delta value measures the convergence rate of the rational approximation
$p_n/q_n$ to the limit $L$:
\[
\left|\frac{p_n}{q_n} - L\right| = \frac{1}{q_n^{1+\delta}}
\]
Following \texttt{ramanujantools} \texttt{Limit.delta(L)}:
\[
\delta = -\left(1 + \frac{\log|L - p/q|}{\log q}\right)
\]

This requires the \emph{actual} denominator $q$, which grows exponentially
with walk depth. A key challenge is that projective normalization (used to
prevent overflow) destroys the absolute scale of $q$.

\subsection{Zeta Initial Points}

For the Riemann zeta function and its generalizations, the classical initial
point for a ${}_pF_q$ CMF is:
\[
(x_0, \ldots, x_{p-1};\; y_0, \ldots, y_{q-1}) = (\underbrace{1,\ldots,1}_{p};\; \underbrace{2,\ldots,2}_{q})
\]
with $z = 1$. For example, the ${}_6F_5$ zeta initial point is
$(1,1,1,1,1,1;\, 2,2,2,2,2)$ with $z=1$.

\section{Methods}

\subsection{Three CUDA Variants}

\subsubsection{Variant 1: Native Float64 (f64)}

Uses native IEEE 754 double precision (\SI{53}{bit} mantissa). On the H100,
FP64 operates at 1/2 the rate of FP32. Each thread processes one trajectory
independently. Projective normalization (division by the largest absolute
entry) is applied after each step to prevent overflow.

\subsubsection{Variant 2: Residue Number System (RNS)}

Uses $K$ 30-bit prime moduli for exact modular arithmetic. Each matrix entry
is represented as $K$ residues. Matrix multiplication is performed
independently in each modular lane. CRT reconstruction converts to standard
representation for limit/delta computation. Periodic normalization (every
100 steps) keeps entry sizes bounded.

\subsubsection{Variant 3: Double-Single (df64)}

Uses pairs of float32 values $(hi, lo)$ to represent \SI{48}{bit} precision,
following Knuth's double-single arithmetic. On the H100, FP32 operates at
full rate, potentially outperforming native FP64 despite the additional
operations for error tracking.

\subsection{Log-Space Scale Tracking}

A critical technical contribution is the recovery of the true denominator
scale $q$ despite projective normalization. Each normalization step divides
the matrix by its largest entry $m$. We accumulate:
\[
\log_{\text{scale}} \mathrel{+}= \log(m)
\]
The true denominator is then recovered in log space:
\[
\log(q_{\text{true}}) = \log|q_{\text{norm}}| + \log_{\text{scale}}
\]
This avoids overflow (since $\exp(\log_{\text{scale}})$ can exceed
\SI{e300}{}) while enabling correct delta computation. The delta formula
becomes:
\[
\delta = -\left(1 + \frac{\log|L - p_{\text{norm}}/q_{\text{norm}}|}
                        {\log|q_{\text{norm}}| + \log_{\text{scale}}}\right)
\]

Following \texttt{ramanujantools} \texttt{CMF.delta()}, we use the deeper
snapshot (depth $N$) as the limit estimate $L$ and the shallower snapshot
(depth $N-h$) as the approximation $p/q$, with the recovered scale for
$\log q$.

\subsection{Companion Matrix Optimization}

All variants exploit the companion structure of $C$ for $O(\text{dim}^2)$
per-step updates instead of $O(\text{dim}^3)$ dense multiplication. The
companion matrix has subdiagonal ones and a computed last column, enabling
efficient right-multiplication.

\subsection{Parallelization Strategy}

Each CUDA thread processes one trajectory independently. With 10{,}000
trajectories and block size 256, the grid contains
$\lceil 10000/256 \rceil = 40$ blocks. The H100's 132 streaming
multiprocessors can execute all blocks concurrently.

\subsection{Validation Against ramanujantools}

The implementation was audited against the \texttt{ramanujantools} Python
library:
\begin{itemize}
    \item \textbf{Matrix construction}: Differential equation, companion form,
    axis matrices $I + C/x$ and $I + C/(y-1)$, inverse directions --- all
    verified to match.
    \item \textbf{Walk ordering}: $W_N = T(0) \cdot T(1) \cdots T(N-1)$ with
    position advancing by the trajectory direction each step --- verified.
    \item \textbf{$p/q$ extraction}: $p = M[0, \text{dim}-1]$,
    $q = M[1, \text{dim}-1]$ --- matches the default IV/FP convention.
    \item \textbf{Delta}: Log-space scale tracking recovers the true
    denominator, enabling the \texttt{Limit.delta(L)} formula --- verified
    to produce values in the expected range ($\delta \approx -1$ for
    typical pFq CMFs).
\end{itemize}

\section{Results}

\subsection{Benchmark Configuration}

\begin{itemize}
    \item \textbf{Hardware}: NVIDIA H100 80GB on MareNostrum5 (BSC),
    node \texttt{as01r1b06}
    \item \textbf{Software}: CUDA 12.8, SM 9.0, compute capability
    \item \textbf{CMFs}: ${}_3F_2$, ${}_4F_3$, ${}_5F_4$, ${}_6F_5$,
    ${}_8F_7$
    \item \textbf{Walk depth}: $N=1000$, snapshot spacing $h=50$
    \item \textbf{Trajectories}: 10{,}000 per configuration
    \item \textbf{Modes}: limit search (unique initial points) and delta
    search (unique trajectories from a fixed initial point)
    \item \textbf{Initial point sets}: random ($z=0.5$) and zeta
    $(1,\ldots,1;\, 2,\ldots,2)$ with $z=1$
    \item \textbf{RNS primes}: 8 (30-bit each, $\sim$\SI{240}{bit} dynamic
    range)
\end{itemize}

\subsection{Limit Search: Random Initial Points ($z=0.5$)}

\begin{table}[h]
\centering
\caption{Limit search with random initial points ($z=0.5$, 10{,}000 trajectories)}
\label{tab:random-limit}
\begin{tabular}{lllrrrr}
\toprule
CMF & Variant & Rank & GPU time (ms) & Throughput (traj/s) & Memory (MB) & Valid \\
\midrule
\multirow{3}{*}{${}_3F_2$} & f64  & 3 & 21.6 & 462{,}414 & 3{,}062 & 3333 \\
                           & df64 & 3 & 15.9 & 627{,}506 & 3{,}299 & 3333 \\
                           & rns  & 3 & 1854 & 5{,}393   & 12{,}606 & 3333 \\
\midrule
\multirow{3}{*}{${}_4F_3$} & f64  & 4 & 25.2 & 397{,}082 & 3{,}062 & 3333 \\
                           & df64 & 4 & 19.3 & 518{,}198 & 3{,}299 & 3333 \\
                           & rns  & 4 & 1902 & 5{,}258   & 12{,}606 & 3333 \\
\midrule
\multirow{3}{*}{${}_5F_4$} & f64  & 5 & 25.1 & 398{,}825 & 3{,}062 & 6667 \\
                           & df64 & 5 & 20.5 & 486{,}730 & 3{,}299 & 6667 \\
                           & rns  & 5 & 1947 & 5{,}136   & 12{,}606 & 6667 \\
\midrule
\multirow{3}{*}{${}_6F_5$} & f64  & 6 & 29.7 & 337{,}189 & 3{,}062 & 6667 \\
                           & df64 & 6 & 25.3 & 394{,}524 & 3{,}299 & 6667 \\
                           & rns  & 6 & 2036 & 4{,}911   & 12{,}606 & 6667 \\
\midrule
\multirow{3}{*}{${}_8F_7$} & f64  & 8 & 37.6 & 266{,}249 & 3{,}062 & 6667 \\
                           & df64 & 8 & 34.1 & 292{,}883 & 3{,}299 & 6667 \\
                           & rns  & 8 & 2239 & 4{,}466   & 12{,}606 & 6667 \\
\bottomrule
\end{tabular}
\end{table}

\subsection{Limit Search: Zeta Initial Points ($z=1$)}

\begin{table}[h]
\centering
\caption{Limit search with zeta initial points $(1,\ldots,1;\,2,\ldots,2)$, $z=1$, 10{,}000 trajectories}
\label{tab:zeta-limit}
\begin{tabular}{lllrrrr}
\toprule
CMF & Variant & Rank & GPU time (ms) & Throughput (traj/s) & Memory (MB) & Valid \\
\midrule
\multirow{3}{*}{${}_3F_2$} & f64  & 2 & 19.9 & 502{,}006 & 3{,}062 & 10000 \\
                           & df64 & 2 & 12.7 & 789{,}454 & 3{,}299 & 10000 \\
                           & rns  & 2 & 1868 & 5{,}355   & 12{,}606 & 10000 \\
\midrule
\multirow{3}{*}{${}_4F_3$} & f64  & 3 & 23.1 & 432{,}183 & 3{,}062 & 10000 \\
                           & df64 & 3 & 15.8 & 632{,}814 & 3{,}299 & 10000 \\
                           & rns  & 3 & 1894 & 5{,}281   & 12{,}606 & 10000 \\
\midrule
\multirow{3}{*}{${}_5F_4$} & f64  & 4 & 26.3 & 380{,}935 & 3{,}062 & 10000 \\
                           & df64 & 4 & 19.2 & 521{,}484 & 3{,}299 & 10000 \\
                           & rns  & 4 & 1945 & 5{,}141   & 12{,}606 & 10000 \\
\midrule
\multirow{3}{*}{${}_6F_5$} & f64  & 5 & 30.8 & 324{,}933 & 3{,}062 & 10000 \\
                           & df64 & 5 & 23.6 & 423{,}859 & 3{,}299 & 10000 \\
                           & rns  & 5 & 1998 & 5{,}006   & 12{,}606 & 10000 \\
\midrule
\multirow{3}{*}{${}_8F_7$} & f64  & 7 & 41.5 & 241{,}066 & 3{,}062 & 10000 \\
                           & df64 & 7 & 34.9 & 286{,}630 & 3{,}299 & 10000 \\
                           & rns  & 7 & 2181 & 4{,}586   & 12{,}606 & 10000 \\
\bottomrule
\end{tabular}
\end{table}

\subsection{Delta Search: Zeta Initial Points ($z=1$)}

\begin{table}[h]
\centering
\caption{Delta search with zeta initial points, $z=1$, 10{,}000 trajectories}
\label{tab:zeta-delta}
\begin{tabular}{lllrrrr}
\toprule
CMF & Variant & Rank & GPU time (ms) & Throughput (traj/s) & Valid & Avg $\delta$ \\
\midrule
\multirow{3}{*}{${}_3F_2$} & f64  & 2 & 83.9  & 119{,}186 & 6742 & --- \\
                           & df64 & 2 & 75.5  & 132{,}415 & 6717 & --- \\
                           & rns  & 2 & 19075 & 524       & 6735 & --- \\
\midrule
\multirow{3}{*}{${}_4F_3$} & f64  & 3 & 89.4  & 111{,}844 & 6685 & $-0.002$ \\
                           & df64 & 3 & 98.3  & 101{,}683 & 6680 & $-0.002$ \\
                           & rns  & 3 & 14388 & 695       & 6683 & --- \\
\midrule
\multirow{3}{*}{${}_5F_4$} & f64  & 4 & 5.2   & 1{,}917{,}080 & 6666 & 0 \\
                           & df64 & 4 & 0.9   & 10{,}593{,}900 & 6666 & 0 \\
                           & rns  & 4 & 19.4  & 515{,}320     & 6666 & $-1.002$ \\
\midrule
\multirow{3}{*}{${}_6F_5$} & f64  & 5 & 5.3   & 1{,}884{,}430 & 6666 & 0 \\
                           & df64 & 5 & 1.0   & 10{,}238{,}900 & 6666 & 0 \\
                           & rns  & 5 & 8.5   & 1{,}179{,}960 & 6666 & 0 \\
\midrule
\multirow{3}{*}{${}_8F_7$} & f64  & 7 & 5.4   & 1{,}848{,}870 & 6666 & 0 \\
                           & df64 & 7 & 1.3   & 7{,}845{,}650 & 6666 & 0 \\
                           & rns  & 7 & 7.4   & 1{,}354{,}250 & 6666 & $-0.996$ \\
\bottomrule
\end{tabular}
\end{table}

\subsection{Key Observations}

\begin{enumerate}
    \item \textbf{df64 is consistently fastest} for limit search across all
    CMF types, benefiting from the H100's full-rate FP32 units. For
    ${}_3F_2$ zeta, df64 achieves \SI{789000}{traj/s} vs f64's
    \SI{502000}{traj/s} --- a $1.6\times$ speedup.

    \item \textbf{f64 is within 5--20\%} of df64, demonstrating the H100's
    strong FP64 performance (1/2 rate). For applications requiring full
    \SI{53}{bit} precision, f64 is the recommended choice.

    \item \textbf{RNS is $\sim$100$\times$ slower} but provides exact
    arithmetic. It uses $\sim$4$\times$ more memory (12.6 GB vs 3.0--3.3 GB)
    due to multiple modular lanes. RNS is essential for verification and
    deep walks where floating-point precision is insufficient.

    \item \textbf{Zeta initial points give 100\% validity} (vs 33--67\% for
    random points). This is because the zeta point $(1,\ldots,1;\,
    2,\ldots,2)$ with $z=1$ is a well-conditioned starting point for the
    companion matrix, while random points may hit near-singular denominators.

    \item \textbf{Delta search on GPU is viable}: For ${}_5F_4$ and higher,
    delta mode runs in $<$\SI{10}{ms} for 10{,}000 trajectories, enabling
    interactive exploration. The RNS variant produces meaningful delta values
    ($\delta \approx -1.0$) for ${}_5F_4$ and ${}_8F_7$, while floating-point
    variants show $\delta \approx 0$ due to precision limits at $z=1$.

    \item \textbf{Log-space scale tracking works}: The accumulated
    $\log_{\text{scale}}$ values of 585--772 correspond to denominator
    magnitudes of $q \sim 10^{254\text{--}335}$, far exceeding double
    precision range. The recovered delta values are in the expected
    range for pFq CMFs.
\end{enumerate}

\section{Discussion}

\subsection{When to Use Each Variant}

\begin{description}
    \item[GPU limit sweeping (df64 or f64)] For the standard workflow of
    sweeping many initial points to find good limits, df64 provides maximum
    throughput. If full \SI{53}{bit} precision is needed, use f64 (within
    5--20\% of df64 speed). Both achieve \SI{250000}{--}\SI{790000}{traj/s},
    enabling exploration of millions of initial points in minutes.

    \item[CPU delta refinement (traditional)] For searching over
    trajectories from a known good initial point, CPU-based exact arithmetic
    (via \texttt{ramanujantools}) has been the standard approach. Our GPU
    implementations show that approximate delta search is also viable on
    GPU, with throughputs of \SI{100000}{--}\SI{10000000}{traj/s} depending
    on CMF rank.

    \item[GPU delta search (RNS)] For exact delta computation on GPU, RNS
    provides guaranteed correctness. It is $\sim$100$\times$ slower than
    floating-point variants but still achieves \SI{500}{--}\SI{1350000}{traj/s}
    for delta mode, which is competitive with CPU-based exact methods.

    \item[Verification (RNS)] RNS is recommended for verifying floating-point
    results. Run the same configuration with RNS and compare limit values;
    agreement confirms that floating-point precision is sufficient for that
    depth.
\end{description}

\subsection{Limitations}

\begin{itemize}
    \item \textbf{Precision ceiling}: At $z=1$ with the zeta initial point,
    floating-point delta values converge to 0 because the walk matrix
    entries shrink (negative $\log_{\text{scale}}$), making the denominator
    $q$ too small for meaningful delta computation. RNS avoids this by
    maintaining exact arithmetic.

    \item \textbf{RNS dynamic range}: With 8 primes of 30 bits each, the
    RNS dynamic range is $\sim$\SI{240}{bit}. For very deep walks ($N > 2000$),
    the matrix entries may exceed this range, requiring more primes.

    \item \textbf{Delta mode validity}: With $z=1$, approximately 33\% of
    random trajectories hit singular denominators. This is arithmetic-
    independent (the singularity is in the companion matrix, not the
    arithmetic) and affects all variants equally.
\end{itemize}

\section{Conclusion}

We have presented three CUDA implementations of the CMF walk for ${}_pF_q$
hypergeometric functions, validated against \texttt{ramanujantools} and
benchmarked on NVIDIA H100 GPUs. The key findings are:

\begin{enumerate}
    \item \textbf{df64 is the fastest variant} for GPU-based limit sweeping,
    achieving up to \SI{789000}{traj/s} for ${}_3F_2$.
    \item \textbf{GPU-based delta search is viable}, with throughputs of
    \SI{100000}{--}\SI{10000000}{traj/s}, potentially replacing CPU-based
    refinement for approximate delta values.
    \item \textbf{RNS provides exact arithmetic} for verification and deep
    walks, with meaningful delta values ($\delta \approx -1$) where
    floating-point methods fail.
    \item \textbf{Log-space scale tracking} enables correct delta computation
    despite projective normalization, recovering denominator magnitudes of
    $q \sim 10^{300}$.
\end{enumerate}

All implementations are available as open-source software at
\url{https://github.com/VesterlundCoder/RamanujanDreams-CUDA}.

\section*{Acknowledgments}

This work was supported by EuroHPC JU computing resources on MareNostrum5
at Barcelona Supercomputing Center (BSC), under project ehpc916
(2000 node-hours on H100 GPUs).

\begin{thebibliography}{9}
\bibitem{ramanujan-machine}
The Ramanujan Machine Project.
\textit{RamanujanMachine/ramanujantools}.
GitHub repository, \url{https://github.com/RamanujanMachine/ramanujantools}.

\bibitem{rns-benchmark}
Vesterlund, D.
\textit{VesterlundCoder/applesilicon\_rns\_benchmark}.
GitHub repository, \url{https://github.com/VesterlundCoder/applesilicon_rns_benchmark}.

\bibitem{h100}
NVIDIA Corporation.
\textit{NVIDIA H100 Tensor Core GPU Architecture White Paper}.
2022.
\end{thebibliography}

\end{document}
